\documentclass[11pt,a4paper]{article}

\usepackage[margin=2.5cm]{geometry}
\usepackage{amsmath,amssymb,amsthm,bm}
\usepackage{booktabs}
\usepackage{hyperref}
\usepackage{enumitem}
\usepackage{listings}
\usepackage{xcolor}
\usepackage{graphicx}

\hypersetup{colorlinks=true,linkcolor=blue,urlcolor=blue,citecolor=blue}

\lstset{
  basicstyle=\ttfamily\small,
  breaklines=true,
  frame=single,
  backgroundcolor=\color{gray!8}
}

\theoremstyle{definition}
\newtheorem{definition}{Definition}
\newtheorem{remark}{Remark}

\title{Harmony Certificate Layer:\\
Mathematical Notes on the Implementation}
\author{Harmony toolbox documentation\\
\texttt{src/Solver/Certificate/}}
\date{\today}

\begin{document}
\maketitle

\begin{abstract}
These notes explain the mathematics behind Harmony's decentralized certificate
layer: passivity, eigenphase, shifted passivity, a Hermitian numerical-range
(DW-style) slice, the five workflows (device gate, PnP library, operating region,
GFM tuning, local vs.\ system), and how they relate to the existing
Leki\'c--Beerten return ratio $H=YZ_{\mathrm{eq}}$.
The implementation uses classical Hermitian passivity and simple phase/shift
screens; it is \emph{inspired by} recent decentralized certificate agendas
but is not a literal port of a specific ETH theorem set.
\end{abstract}

\tableofcontents
\newpage

%==============================================================================
\section{Setup and notation}
%==============================================================================

\subsection{Device admittance}

At a fixed operating point, after equilibrium and linearization
(\texttt{solveEquilibrium}, \texttt{computeABCD}), each element exposes a
frequency-domain admittance
\begin{equation}
  Y(j\omega)\;\in\;\mathbb{C}^{n\times n},
  \qquad \omega = 2\pi f.
\end{equation}
In code this is \texttt{Element::compute\_y\_parameters(f)} converted to an
Eigen matrix.

For an MMC, Harmony typically builds a multiport map with DC and AC $dq$
ports. The certificate layer can use:
\begin{itemize}[leftmargin=*]
  \item the full matrix, or
  \item the \textbf{AC $dq$ block} (default): if $Y\in\mathbb{C}^{3\times 3}$
  with ordering (DC, $d$, $q$), take
  \begin{equation}
    Y_{\mathrm{ac}}(j\omega)
    \;=\;
    Y(j\omega)\big|_{[2:3],[2:3]}
    \;\in\;\mathbb{C}^{2\times 2}.
  \end{equation}
  (\texttt{extractAcDqBlock}; 0-based indices $1\!:\!2$ in Eigen.)
\end{itemize}

Unless stated otherwise, write $Y(j\omega)$ for the matrix actually tested
(full or AC block).

\subsection{Hermitian part}

\begin{definition}[Hermitian part]
\[
  \operatorname{Her}(Y)
  \;:=\;
  \tfrac12\bigl(Y+Y^*\bigr),
  \qquad
  Y^*=\overline{Y}^\top.
\]
\end{definition}
$\operatorname{Her}(Y)$ is always Hermitian, so its eigenvalues are real.

\subsection{Frequency grid}

Sweeps use a log-spaced grid
\begin{equation}
  f_k
  =
  10^{\,\log_{10} f_{\min}
  + \frac{k}{N-1}(\log_{10} f_{\max}-\log_{10} f_{\min})},
  \quad k=0,\ldots,N-1,
\end{equation}
with defaults typically $f_{\min}=1\,\mathrm{Hz}$, $f_{\max}=1000\,\mathrm{Hz}$,
$N\approx 60$--$120$ (\texttt{CertificateSpec}).

%==============================================================================
\section{Core metrics (what the code computes)}
%==============================================================================

\subsection{Passivity index}

\begin{definition}[Passivity index]
\[
  \nu\bigl(Y(j\omega)\bigr)
  \;:=\;
  \lambda_{\min}\!\bigl(\operatorname{Her} Y(j\omega)\bigr).
\]
\end{definition}
Code: \texttt{passivityIndex(Y)}.

\begin{remark}[Passivity test]
$Y(j\omega)$ is (strictly) passive at $\omega$ if $\nu(Y)>0$, and
passive in the weak sense if $\nu(Y)\ge 0$. Harmony uses
\[
  \nu(Y)\;\ge\;\varepsilon
\]
with default $\varepsilon=0$ (\texttt{passivity\_eps}), i.e.\ positive
\emph{semi}definite Hermitian part.
\end{remark}

\textbf{Interpretation.}
For currents/voltages related by $I=YV$, passivity of $Y$ is a classical
dissipativity condition that composes under interconnection of passive
systems. It is a \emph{sufficient} local screen for plug-and-play style
arguments, not necessary for closed-loop stability.

Worst-case over the grid:
\begin{equation}
  \nu_{\mathrm{worst}}
  \;=\;
  \min_{k}\,\nu\bigl(Y(j\omega_k)\bigr).
\end{equation}

\subsection{Eigenphase indicator}

Let $\lambda_i(Y)$ be the eigenvalues of $Y$ (not of $\operatorname{Her} Y$).
\begin{definition}[Max eigenphase]
\[
  \varphi(Y)
  \;:=\;
  \max_{i:\,\lambda_i\neq 0}
  \bigl|\arg\lambda_i(Y)\bigr|
  \qquad\text{(radians)},
\]
reported in degrees as $\varphi_{\deg}=(180/\pi)\,\varphi$.
\end{definition}
Code: \texttt{maxEigenPhaseDeg(Y)}.

Pass if
\begin{equation}
  \varphi_{\deg}(Y)
  \;\le\;
  \varphi_{\mathrm{lim}}
\end{equation}
(default $\varphi_{\mathrm{lim}}=90^\circ$).

\begin{remark}
This is a \textbf{simple MIMO phase proxy}, not a graph phase margin and not
a full sectorial / DW-shell certificate. Eigenvalues near the negative real
axis give $\varphi\approx 180^\circ$ and fail a $90^\circ$ limit even if
$\operatorname{Her} Y\succeq 0$ (seen on some passive branch stamps).
\end{remark}

Worst-case:
\begin{equation}
  \varphi_{\mathrm{worst}}
  \;=\;
  \max_{k}\,\varphi_{\deg}\bigl(Y(j\omega_k)\bigr).
\end{equation}

\subsection{Shifted passivity (loop-shift flavour)}

\begin{definition}[Shifted passivity]
For a real shift $\delta$,
\[
  \nu_\delta(Y)
  \;:=\;
  \lambda_{\min}\!\bigl(\operatorname{Her}(Y+\delta I)\bigr).
\]
\end{definition}
Code: \texttt{shiftedPassivityIndex(Y, delta)}.

This is a \textbf{constant} feedthrough shift: if $\delta>0$ is large enough,
$\nu_\delta$ can become positive even when $\nu(Y)<0$. It mimics the
\emph{idea} of loop shifting used to handle non-passive devices/lines, but
it is not dynamic loop-shifting of the full GFM literature.

\subsection{Hermitian DW-style slice}

For Hermitian $H=\operatorname{Her} Y$, the numerical range is the real
interval
\begin{equation}
  W(H)=\bigl[\lambda_{\min}(H),\,\lambda_{\max}(H)\bigr]\subset\mathbb{R}.
\end{equation}
Harmony's DW margin is defined as
\begin{equation}
  m_{\mathrm{DW}}(Y)
  \;:=\;
  \lambda_{\min}(\operatorname{Her} Y)
  \;=\;
  \nu(Y).
\end{equation}
Code: \texttt{dwHermitianMargin} (same number as passivity).

The plot \texttt{plot\_certificate\_dw\_slice} draws an ellipse in the plane
with semi-axes $|\lambda_{\min}|$, $|\lambda_{\max}|$ as a visual stand-in for
a numerical-range boundary. This is \textbf{not} a full Davis--Wielandt shell
or SRG certificate.

\subsection{Extra diagnostic}

The sweep also stores
\begin{equation}
  \min_i \operatorname{Re}\lambda_i\bigl(Y(j\omega)\bigr),
\end{equation}
useful for plotting but not used as a hard gate criterion by default.

%==============================================================================
\section{Certificate verdict}
%==============================================================================

Given a completed sweep and a \texttt{CertificateSpec},
\begin{align}
  \mathrm{pass\_passivity}
  &\iff \nu_{\mathrm{worst}}\ge\varepsilon,\\
  \mathrm{pass\_phase}
  &\iff \varphi_{\mathrm{worst}}\le\varphi_{\mathrm{lim}},\\
  \mathrm{pass\_shifted}
  &\iff \nu_{\delta,\mathrm{worst}}\ge\varepsilon,\\
  \mathrm{pass\_dw}
  &\iff m_{\mathrm{DW,worst}}\ge\varepsilon.
\end{align}
Enabled checks are selected by Boolean flags
\texttt{require\_passivity}, \texttt{require\_phase}, \texttt{require\_shifted}.

%==============================================================================
\section{System return ratio $H=YZ_{\mathrm{eq}}$}
%==============================================================================

Harmony's \texttt{StabilityEstimate} already computes the converter-facing
return ratio (Leki\'c--Beerten style)
\begin{equation}
  H(j\omega)
  \;=\;
  Y_{\mathrm{conv}}(j\omega)\,Z_{\mathrm{eq}}(j\omega),
\end{equation}
where $Z_{\mathrm{eq}}$ is the network equivalent seen at the chosen cut
(\texttt{compute\_transfer\_function}).

For the \textbf{system-side} certificate sweep, Harmony evaluates passivity of
\begin{equation}
  I+H(j\omega)
\end{equation}
via
\begin{equation}
  \nu_{\mathrm{sys}}(\omega)
  \;:=\;
  \lambda_{\min}\!\bigl(\operatorname{Her}\bigl(I+H(j\omega)\bigr)\bigr).
\end{equation}
Phase is still reported on $H$ itself:
\begin{equation}
  \varphi_{\mathrm{sys}}(\omega)
  \;=\;
  \varphi_{\deg}\bigl(H(j\omega)\bigr).
\end{equation}

\begin{remark}[Why $I+H$?]
Nyquist-type interconnection stability for the feedback $H$ is related to
whether $-1$ is avoided by the eigenvalues of $H$. Looking at
$\operatorname{Her}(I+H)\succ 0$ is a convenient \emph{passivity-style proxy}
on the closed-loop return operator, not a complete MIMO Nyquist theorem.
Use it as a benchmark alongside Bode/Nyquist of $H$, not as a replacement.
\end{remark}

%==============================================================================
\section{The five workflows}
%==============================================================================

\subsection{Workflow 1 --- Device gate}

\textbf{API:} \texttt{evaluateDeviceGate(elem, id, spec)}.

Compute the device sweep on $Y$, then
\begin{equation}
  \mathrm{ALLOW}
  \iff
  \bigl(\neg R_\nu \lor \mathrm{pass\_passivity}\bigr)
  \;\land\;
  \bigl(\neg R_\varphi \lor \mathrm{pass\_phase}\bigr)
  \;\land\;
  \bigl(\neg R_\delta \lor \mathrm{pass\_shifted}\bigr),
\end{equation}
where $R_\nu,R_\varphi,R_\delta$ are the \texttt{require\_*} flags.

\textbf{Product meaning:} a local allow/block before interconnecting a device
into an unknown or changing network (PnP gate).

\subsection{Workflow 2 --- PnP library}

\textbf{API:} \texttt{PnPCertificateLibrary}.

Store records
\[
  \bigl(\mathrm{id},\;\mathrm{type},\;\mathrm{allow},\;
  \nu_{\mathrm{worst}},\;\varphi_{\mathrm{worst}},\;\ldots\bigr)
\]
in memory and persist as JSON/CSV. Lookup: \texttt{isAllowed(id)}.

No new mathematics --- it is a catalog of gate results for fleet screening.

\subsection{Workflow 3 --- Operating $(P,Q)$ region}

\textbf{API:} \texttt{certifyOperatingRegion}, \texttt{certifyOperatingPoints}.

On a rectangular grid
\begin{equation}
  P\in[P_{\min},P_{\max}],\quad
  Q\in[Q_{\min},Q_{\max}],
\end{equation}
for each sample $(P_i,Q_j)$:
\begin{enumerate}[leftmargin=*]
  \item \texttt{update\_MMC}$(V_m,\theta,P,Q,V_{\mathrm{dc}},P_{\mathrm{dc}})$,
  \item re-solve equilibrium and rebuild ABCD/$Y$,
  \item run the device certificate sweep,
  \item mark the point certified iff enabled checks pass.
\end{enumerate}
Optional post-OPF check: also evaluate the solved OPF injection $(P^\star,Q^\star)$
and report whether that setpoint is certified.

Margin stored per point is typically $\nu_{\mathrm{worst}}$ at that OP.

\subsection{Workflow 4 --- GFM droop tuning}

GFM outer loop (Harmony MMC mode) uses droops $(K_P,K_Q)$ stored as the PI
gains of the \texttt{gfm} controller
(\texttt{getGfmDroops}/\texttt{setGfmDroops}).

On a search grid over $(K_P,K_Q)$:
\begin{enumerate}[leftmargin=*]
  \item set droops, re-equilibrate, rebuild $Y$,
  \item evaluate the device certificate,
  \item among feasible points, keep the one with largest $\nu_{\mathrm{worst}}$.
\end{enumerate}
If none feasible, restore nominal droops.

This is a \textbf{numerical tuning assistant}, not a closed-form parametric
certificate set from GFM theory papers.

\subsection{Workflow 5 --- Local vs.\ system}

\textbf{API:} \texttt{compareLocalVsSystem}.

Compute
\begin{align}
  \mathrm{local\_pass}
  &\text{ from device }Y\text{ (enabled checks)},\\
  \mathrm{system\_pass}
  &\iff \nu_{\mathrm{sys,worst}}\ge\varepsilon
  \quad\bigl(\operatorname{Her}(I+H)\bigr).
\end{align}
Report consistency:
\begin{equation}
  \mathrm{consistent}
  \iff
  (\mathrm{local\_pass}=\mathrm{system\_pass}).
\end{equation}

The interesting empirical case for GFM-MMC is often
\begin{equation}
  \mathrm{local\_pass}=\mathrm{false},
  \qquad
  \mathrm{system\_pass}=\mathrm{true},
\end{equation}
i.e.\ PnP/passivity rejects a device that still looks acceptable under the
network-aware $H$ test --- the conservatism story for a methods/application
paper.

%==============================================================================
\section{MMC multiport picture (for AC--DC reading)}
%==============================================================================

Conceptually an MMC at small signal is a multiport admittance relating
perturbations of DC and AC $dq$ voltages/currents. Harmony's $3\times 3$ $Y$
can be partitioned as
\begin{equation}
  Y
  =
  \begin{bmatrix}
    Y_{\mathrm{dd}} & Y_{\mathrm{da}} \\
    Y_{\mathrm{ad}} & Y_{\mathrm{aa}}
  \end{bmatrix},
  \qquad
  Y_{\mathrm{aa}}\in\mathbb{C}^{2\times 2}.
\end{equation}
Today's default gate uses $Y_{\mathrm{aa}}$. A natural AC--DC extension is to
certify ports separately,
\begin{equation}
  \nu_{\mathrm{ac}}=\nu(Y_{\mathrm{aa}}),\qquad
  \nu_{\mathrm{dc}}=\nu(Y_{\mathrm{dd}}),\qquad
  \nu_{\mathrm{full}}=\nu(Y),
\end{equation}
and/or compare local AC/DC screens to $H$ cuts on AC and DC sides.

%==============================================================================
\section{GFM outer-loop reminder (context for tuning)}
%==============================================================================}

In Harmony's GFM mode, filtered powers $(P_f,Q_f)$ and angle $\theta$ evolve
with droops $(K_P,K_Q)$ and an internal voltage magnitude reference. The
linearized plant+control yields $Y(j\omega)$ used by the certificates.
Changing $(K_P,K_Q)$ changes $Y$ and therefore $\nu$ and $\varphi$; the
tuning workflow searches that map numerically.

%==============================================================================
\section{How to run and what files mean}
%==============================================================================

\begin{lstlisting}
Harmony --cpp certificate_figures
Harmony --cpp certificate_figures --no-plot
\end{lstlisting}

Typical outputs under \texttt{files/}:
\begin{center}
\begin{tabular}{@{}ll@{}}
\toprule
File & Content \\
\midrule
\texttt{certificate\_device\_sweep.csv} & $f,\nu,\nu_\delta,\varphi,\ldots$ \\
\texttt{certificate\_system\_H\_sweep.csv} & same for $I+H$ / $H$ \\
\texttt{certificate\_operating\_region.csv} & $(P,Q)$, certified, margin \\
\texttt{pnp\_certificate\_library.json} & gate catalog \\
\texttt{certificate\_tune\_before/after.csv} & tuning sweeps \\
\bottomrule
\end{tabular}
\end{center}

Source map:
\begin{center}
\begin{tabular}{@{}ll@{}}
\toprule
Header & Role \\
\midrule
\texttt{Certificate\_spec.h} & thresholds / grids \\
\texttt{Stability\_certificate.h} & metrics + plots \\
\texttt{Device\_gate.h} & workflow 1 \\
\texttt{PnP\_library.h} & workflow 2 \\
\texttt{Operating\_region.h} & workflow 3 \\
\texttt{Tuning\_assistant.h} & workflow 4 \\
\texttt{Local\_vs\_system.h} & workflow 5 \\
\bottomrule
\end{tabular}
\end{center}

%==============================================================================
\section{What this is / is not (learning checklist)}
%==============================================================================

\begin{center}
\begin{tabular}{@{}p{0.46\textwidth}p{0.46\textwidth}@{}}
\toprule
\textbf{Implemented} & \textbf{Not implemented} \\
\midrule
$\lambda_{\min}(\operatorname{Her} Y)$ passivity &
Full DW-shell / SRG geometry \\
Simple eigenphase of $Y$ &
Graph phase-margin certificates \\
Constant shift $\delta$ &
Dynamic loop-shifting GFM theorems \\
Local vs $\operatorname{Her}(I+H)$ &
Proof that local $\Rightarrow$ system \\
$(P,Q)$ re-linearization sweep &
Analytic OP family certificates \\
GFM droop grid search &
Closed-form allowable gain sets \\
\bottomrule
\end{tabular}
\end{center}

%==============================================================================
\section*{Suggested reading path}
%==============================================================================

\begin{enumerate}[leftmargin=*]
  \item Re-derive $\nu(Y)=\lambda_{\min}(\operatorname{Her} Y)$ on a $2\times 2$
  resistive $Y$ by hand; confirm \texttt{passivityIndex}.
  \item Run \texttt{certificate\_figures --no-plot}; open the CSV sweeps.
  \item Compare device $\nu(\omega)$ to system $\nu_{\mathrm{sys}}(\omega)$
  and note FAIL/PASS disagreement on GFM.
  \item Only then read D\"orfler-line PnP / DW-shell papers as
  \emph{extensions} beyond this toolbox layer.
\end{enumerate}

\end{document}
